 \begin{titlepage}

\begin{center}

% %% Extra whitespace at the top.
% \vspace*{2\bigskipamount}

% %% Print the title.
% {\makeatletter
% \titlestyle\bfseries\LARGE\@title
% \makeatother}

% %% Print the optional subtitle.
% {\makeatletter
% \ifx\@subtitle\undefined\else
%     \bigskip
%     \titlefont\titleshape\Large\@subtitle
% \fi
% \makeatother}

% \end{center}

% \cleardoublepage
% \thispagestyle{empty}

% \begin{center}

%% The following lines repeat the previous page exactly.

\vspace*{2\bigskipamount}

%% Print the title.
{\makeatletter
\titlestyle\bfseries\LARGE\@title
\makeatother}

%% Print the optional subtitle.
{\makeatletter
\ifx\@subtitle\undefined\else
    \bigskip
    \titlefont\titleshape\Large\@subtitle
\fi
\makeatother}

%% Uncomment the following lines to insert a vertically centered picture into
%% the title page.
%\vfill
%\includegraphics{title}
% \vfill

%% Apart from the names and dates, the following text is dictated by the
%% promotieregelement.
\bigskip
\bigskip
\bigskip
\bigskip
\bigskip
\bigskip

{\Large\titlefont Proefschrift}

\bigskip
\bigskip

ter verkrijging van

de graad van Doctor aan de Universiteit Leiden,

op gezag van Rector Magnificus prof.mr. C.J.J.M. Stolker,

volgens besluit van het College voor Promoties

te verdedigen op dinsdag 3 april 2018

klokke 15.00 uur

\bigskip
\bigskip

door

\bigskip
\bigskip

%% Print the full name of the author.
\makeatletter
{\Large\titlefont\bfseries\@firstname\ {\titleshape\@lastname}}
\makeatother

\bigskip
\bigskip


geboren te Angul (India) \\
in 1991

%% Extra whitespace at the bottom.
\vspace*{2\bigskipamount}

\end{center}

\clearpage
\thispagestyle{empty}

%% The following line is dictated by the promotieregelement.
\noindent Promotores:

%% List the promotors (supervisors).
\medskip\noindent
\begin{tabular}{l}
    Prof.\ dr.\ M.A.G.J.\ Orrit \\
    Prof.\ dr.\ G.W.\ Canters
\end{tabular}

% % List the (optional) copromotor.
% \medskip
% \noindent Copromotor: Prof.\ dr.\ T.J\ Aartsma

\medskip
\noindent Promotiecommissie:

%% List the committee members, starting with the Rector Magnificus and the
%% promotor(s) and ending with the reserve members.
\medskip\noindent
\begin{tabular}{ll}
    Dr.\ K.\ Blank (Max Planck Institute of Colloids and Interfaces) \\
    Dr.\ P.\ Zijlstra (Technische Universiteit Eindhoven) \\
    Prof.\ dr.\ E.R.\ Eliel \\
    Prof.\ dr.\ E.J.J.\ Groenen \\
    Prof.\ dr.\ A.\ Kros \\
    Prof.\ dr.\ ir.\ S.J.T. van Noort \\
\end{tabular}

%% Include the following disclaimer for committee members who have contributed
%% to this dissertation. Its formulation is again dictated by the
%% promotieregelement.
% \medskip
% \noindent Prof.\ dr.\ ir.\ J.\ Doe heeft als begeleider in belangrijke mate aan de totstandkoming van het proefschrift bijgedragen.

%% Here you can include the logos of any institute that contributed financially
%% to this dissertation.
\vfill
\begin{center}
    \includegraphics[height=0.8in]{frontback/logos/leiden}
    \hspace{2em}
    \includegraphics[height=0.8in]{frontback/logos/casimir} \\
%    \hspace{2em}
%    \includegraphics[height=0.5in]{title/logos/nwo}
\end{center}
\vfill

\vspace{2\bigskipamount}

\noindent Copyright \textcopyright\ 2018 by B.~Pradhan

%% Uncomment the following lines if this dissertation is part of the Casimir PhD
%% Series, or a similar research school.
%\medskip
%\noindent Casimir PhD Series, Delft-Leiden 2013-01

\medskip
\noindent Casimir PhD series, Delft-Leiden 2018-08 \\
\noindent ISBN 978-90-8593-340-3

\medskip
\noindent An electronic version of this dissertation is available at \\
\url{http://repository.tudelft.nl/}.

\medskip
\noindent The work described in this thesis is supported by the NanoFront consortium, a program of the Netherlands Organisation for Scientific Research (NWO) that is funded by the Dutch Ministry of Education, Culture and Science (OCW).

\medskip
\noindent The cover shows cartoon structures of an azurin and a DNA around a gold nanorod. The background is an artistic picture of a night sky where the bright spots are analogs of stars and planets. Observing single molecules by optical microscopy can be imagined as if you are watching blinking stars in the night sky.  

\end{titlepage}

